\section{Einleitung und Zielsetzung}

\subsection{Ausgangslage}

Öffentliche Insider-Trade-Daten können über API-Schnittstellen bezogen und als Grundlage einer analytischen Datenanwendung genutzt werden \parencite{fmp_docs}. Für eine belastbare Auswertung reicht der direkte Zugriff auf die API-Antworten jedoch nicht aus. Die gelieferten Rohdaten müssen zunächst normalisiert, validiert und fachlich gefiltert werden, bevor sie als konsistente Analysegrundlage verwendet werden können.

Mercator ist vor diesem Hintergrund als analytische Datenanwendung einzuordnen. Analytische Anwendungen unterscheiden sich von transaktionalen Systemen dadurch, dass nicht die operative Verarbeitung einzelner Geschäftsvorfälle im Vordergrund steht, sondern die Auswertung, Aggregation und Exploration von Daten \parencite{ibm_olap_oltp}. Kipf beschreibt den steigenden Bedarf an analytischer Datenverarbeitung ebenfalls als Folge wachsender Datenmengen und zunehmender Datenverfügbarkeit \parencite{kipf_2020_analytical_database_systems}. Für Mercator bedeutet dies, dass die Anwendung nicht primär einzelne Transaktionen verarbeitet, sondern Insider-Trade-Daten strukturiert aufbereitet, speichert, bewertet und über eine Oberfläche analysierbar macht.

\subsection{Problemstellung}

Insider-Trade-Daten besitzen einen fachlichen Analysewert, sind jedoch nicht ohne weitere Verarbeitung interpretierbar. Cohen, Malloy und Pomorski zeigen, dass Insider-Trades differenziert betrachtet werden müssen, da insbesondere opportunistische Transaktionen einen höheren Informationsgehalt besitzen können als routinemäßige Transaktionen \parencite{cohen_malloy_pomorski_2012}. Daraus folgt für Mercator, dass nicht jeder geladene Datensatz automatisch als relevant behandelt werden darf.

Die zentrale technische Herausforderung besteht darin, rohe API-Daten in eine nachvollziehbare Verarbeitungskette zu überführen. Dafür müssen ungültige oder unvollständige Datensätze erkannt, fachliche Filterregeln angewendet und relevante Trades gezielt angereichert werden. Weiterführende API-Abfragen sollen nur dann ausgelöst werden, wenn ein Datensatz zuvor lokale Mindestkriterien erfüllt. Dadurch wird die externe API geschont und die Verarbeitung bleibt reproduzierbar. Zusätzlich müssen unterschiedliche Speicheranforderungen berücksichtigt werden: Rohdaten sollen nachvollziehbar erhalten bleiben, während bereinigte Daten in einer strukturierten Form für Analyse und Visualisierung verfügbar sein müssen.

\subsection{Ziel der Arbeit}

Ziel der Arbeit ist die Entwicklung und Dokumentation einer prototypischen analytischen Datenanwendung zur Verarbeitung, Bewertung und Visualisierung öffentlicher Insider-Trade-Daten. Die Anwendung soll Daten aus der Financial-Modeling-Prep-API laden, lokal validieren, fachlich filtern, mit Unternehmens- und Marktdaten anreichern und anschließend in einer interaktiven Streamlit-Oberfläche darstellen.

Die technische Umsetzung kombiniert MongoDB und MySQL. MongoDB dient als Raw Store für rohe beziehungsweise flexible API-Payloads. MySQL dient als Clean Store für bereinigte, normalisierte und analysierbare Daten. Ergänzend wird ein regelbasiertes Scoring umgesetzt, das die Priorisierung von Datensätzen transparent und nachvollziehbar macht. Der Schwerpunkt liegt damit auf Datenmodellierung, Datenverarbeitung, Speicherarchitektur und Anwendungsgestaltung.

\subsection{Forschungs- und Projektfrage}

Die Arbeit orientiert sich an folgender zentraler Projektfrage:

\begin{quote}
Wie kann eine analytische Datenanwendung aufgebaut werden, die öffentliche Insider-Trade-Daten nachvollziehbar verarbeitet, speichert, bewertet und interaktiv visualisiert?
\end{quote}

Zur Beantwortung dieser Frage werden die wesentlichen Verarbeitungsschritte von der Datenaufnahme über Validierung, Speicherung und Anreicherung bis zur Visualisierung beschrieben. Dabei wird insbesondere untersucht, wie eine Raw/Clean-Trennung, lokale Gate-Regeln und ein transparentes Scoring innerhalb einer kompakten Projektarchitektur umgesetzt werden können.

\subsection{Abgrenzung}

Mercator ist keine Anlageberatung und erzeugt keine Kauf- oder Verkaufsempfehlungen. Das Scoring ist keine Prognose über zukünftige Kursentwicklungen, sondern eine regelbasierte Heuristik zur strukturierten Priorisierung von Analyseobjekten. Die Anwendung ist außerdem kein produktives Handelssystem und adressiert keine Echtzeit-Entscheidungen. Der Fokus dieser Arbeit liegt auf der technischen Umsetzung einer Datenanwendung im Kontext des Moduls Datenbanken 2. Finanztheoretische Bewertung, Marktvoraussagen und die Einbindung weiterer externer Datenquellen außerhalb der verwendeten API-Schnittstellen sind nicht Gegenstand dieser Arbeit.